6.1 systematische Fehler
Kalibrierungsfehler:
Eine fehlerhafte Kalibrierung der Messgeräte verursacht systematische Messfehler. Wurde die Kalibrierung nicht ordnungsgemäß vorgenommen, weichen die ermittelten Messwerte von den physikalisch tatsächlichen Werten ab. Um diese Abweichung so gering wie möglich zu halten, sind die Messgeräte regelmäßig und gewissenhaft zu kalibrieren – am besten unmittelbar vor jeder Messreihe.
Temperaturschwankungen:
Temperaturschwankungen verändern die elektrischen Eigenschaften der Bauelemente, vor allem die Kapazität des Kondensators sowie die Kennwerte der gesamten Schaltung. Eine konstante Umgebungstemperatur während der Messung oder der Einsatz von Temperaturkompensationsmaßnahmen tragen dazu bei, diese Fehlerquelle einzudämmen.
Parasitische Kapazitäten:
Parasitäre Kapazitäten, verursacht durch metallische Gegenstände in der Nähe oder die Verdrahtung der Schaltung, führen zu verfälschten Messergebnissen. Durch einen sorgfältig ausgelegten Versuchsaufbau sowie abgeschirmte Leitungen und Bauteile lassen sich diese Störeffekte deutlich verringern.
6.2 Zufällige Fehler
Rauschen und Störungen:
Elektronisches Rauschen sowie elektromagnetische Störungen beeinträchtigen die Signalgüte und verursachen Ungenauigkeiten bei den Messungen. Um diesen Einfluss einzudämmen, sind gut abgeschirmte Messgeräte einzusetzen und der Versuchsaufbau in einer störungsarmen Umgebung aufzubauen. Zusätzlich lassen sich die Messdaten mit digitalen Filtern bearbeiten, um Störanteile herauszufiltern.
Messunsicherheiten durch digitale Zählung:
Bei der Frequenzmessung mittels digitaler Zählschaltungen treten Zählfehler auf, welche die Messgenauigkeit herabsetzen. Eine Verlängerung der Messzeit sowie eine Erhöhung der Abtastrate senken die Messunsicherheit und ermöglichen genauere Messergebnisse. Aus dem verwendeten Skript ergeben sich konkrete Optimierungsfragen für die Frequenzbestimmung. Beispielsweise ist zu klären, wie viele Messwerte erfasst werden müssen, um eine vorgegebene Messunsicherheit einzuhalten.
